% Fact sheets for rewriting the paper in your own words. NOT part of the paper.
% Numbers come from numbers.tex (same as the paper), so they always match.
% Build: tectonic factsheets.tex   ->  factsheets.pdf
\documentclass[11pt]{article}
\usepackage[margin=1in]{geometry}
\usepackage{enumitem}
\usepackage{xcolor}
\setlist{itemsep=2pt, topsep=3pt}
\newcommand{\todo}[1]{\textcolor{red}{[TODO: #1]}}
% Generated by src/make_numbers.py from the result files. Do not edit by hand.
\newcommand{\actbuystocks}{2}
\newcommand{\activeA}{69}
\newcommand{\activeL}{39}
\newcommand{\activeM}{49}
\newcommand{\activeN}{276}
\newcommand{\activeP}{53}
\newcommand{\aucbsA}{0.67}
\newcommand{\aucbsL}{0.64}
\newcommand{\aucbsM}{0.66}
\newcommand{\aucbsN}{0.63}
\newcommand{\aucbsP}{0.63}
\newcommand{\aucchronosM}{0.57}
\newcommand{\aucfintextM}{0.52}
\newcommand{\auclgbM}{0.66}
\newcommand{\auctimesfmM}{0.58}
\newcommand{\baseA}{59.4}
\newcommand{\baseL}{49.7}
\newcommand{\baseM}{56.7}
\newcommand{\baseN}{53.7}
\newcommand{\baseP}{54.0}
\newcommand{\basepooledA}{56.9}
\newcommand{\basepooledL}{45.8}
\newcommand{\basepooledM}{56.4}
\newcommand{\basepooledN}{52.6}
\newcommand{\basepooledP}{52.4}
\newcommand{\bigM}{+2.4}
\newcommand{\bigerrM}{5}
\newcommand{\bighiM}{4.4}
\newcommand{\bigloM}{0.3}
\newcommand{\bigmax}{+4.3}
\newcommand{\bigmin}{\ensuremath{-}1.7}
\newcommand{\bsaucM}{0.66}
\newcommand{\bsbaseA}{69.7}
\newcommand{\bsbaseL}{65.0}
\newcommand{\bsbaseM}{64.7}
\newcommand{\bsgainA}{+0.7}
\newcommand{\bsgainL}{+2.1}
\newcommand{\bsgainM}{+2.7}
\newcommand{\bslgbA}{70.4}
\newcommand{\bslgbL}{67.1}
\newcommand{\bslgbM}{67.4}
\newcommand{\chronosDM}{56.5}
\newcommand{\chronosM}{56.5}
\newcommand{\chronosholdM}{0.0}
\newcommand{\chronospooledM}{58.3}
\newcommand{\cpbuysA}{\ensuremath{-}0.14}
\newcommand{\cpbuysL}{+0.46}
\newcommand{\cpbuysM}{+0.20}
\newcommand{\cpbuystA}{\ensuremath{-}0.7}
\newcommand{\cpbuystL}{1.8}
\newcommand{\cpbuystM}{1.0}
\newcommand{\cpsellsA}{+0.01}
\newcommand{\cpsellsL}{\ensuremath{-}0.07}
\newcommand{\cpsellsM}{+0.31}
\newcommand{\cpsellstA}{0.0}
\newcommand{\cpsellstL}{\ensuremath{-}0.2}
\newcommand{\cpsellstM}{1.8}
\newcommand{\cusipchangen}{6}
\newcommand{\cusipchangepct}{0.01}
\newcommand{\differrA}{+3.7}
\newcommand{\differrM}{+2.9}
\newcommand{\differrNM}{\ensuremath{-}3.3}
\newcommand{\differrNP}{+1.3}
\newcommand{\differrhiA}{10.7}
\newcommand{\differrhiM}{10.2}
\newcommand{\differrhiNM}{2.1}
\newcommand{\differrhiNP}{5.8}
\newcommand{\differrhinoetfM}{11.4}
\newcommand{\differrloA}{\ensuremath{-}3.7}
\newcommand{\differrloM}{\ensuremath{-}4.7}
\newcommand{\differrloNM}{\ensuremath{-}8.6}
\newcommand{\differrloNP}{\ensuremath{-}3.3}
\newcommand{\differrlonoetfM}{\ensuremath{-}3.9}
\newcommand{\differrnoetfM}{+4.1}
\newcommand{\diffgainA}{+2.7}
\newcommand{\diffgainM}{+2.1}
\newcommand{\diffgainNM}{\ensuremath{-}1.2}
\newcommand{\diffgainNP}{+0.6}
\newcommand{\diffgainhiA}{6.5}
\newcommand{\diffgainhiM}{6.0}
\newcommand{\diffgainhiNM}{1.2}
\newcommand{\diffgainhiNP}{2.7}
\newcommand{\diffgainhinoetfM}{6.8}
\newcommand{\diffgainloA}{\ensuremath{-}1.1}
\newcommand{\diffgainloM}{\ensuremath{-}1.9}
\newcommand{\diffgainloNM}{\ensuremath{-}4.0}
\newcommand{\diffgainloNP}{\ensuremath{-}1.7}
\newcommand{\diffgainlonoetfM}{\ensuremath{-}1.6}
\newcommand{\diffgainnoetfM}{+2.7}
\newcommand{\diffsamplebuys}{+0.11}
\newcommand{\diffsamplebuyst}{1.9}
\newcommand{\diffsamplesells}{+0.15}
\newcommand{\diffsamplesellst}{2.5}
\newcommand{\dmlltmin}{8.0}
\newcommand{\dmpchronos}{.033}
\newcommand{\dmpfintext}{.11}
\newcommand{\dmposM}{18}
\newcommand{\dmposchronos}{6}
\newcommand{\dmposmin}{16}
\newcommand{\dmpostimesfm}{1}
\newcommand{\dmqM}{18}
\newcommand{\dmtM}{7.4}
\newcommand{\dmtchronos}{\ensuremath{-}2.3}
\newcommand{\dmtfintext}{\ensuremath{-}1.7}
\newcommand{\dmtminall}{5.3}
\newcommand{\dmttimesfm}{\ensuremath{-}5.6}
\newcommand{\dwgain}{+2.8}
\newcommand{\dwhi}{5.2}
\newcommand{\dwlo}{0.9}
\newcommand{\earlyA}{+2.0}
\newcommand{\earlyL}{+6.2}
\newcommand{\earlyM}{+0.6}
\newcommand{\earlyerrA}{5}
\newcommand{\earlyerrL}{12}
\newcommand{\earlyerrM}{1}
\newcommand{\earlyhiA}{3.5}
\newcommand{\earlyhiL}{11.7}
\newcommand{\earlyhiM}{2.0}
\newcommand{\earlyloA}{0.6}
\newcommand{\earlyloL}{2.0}
\newcommand{\earlyloM}{\ensuremath{-}1.0}
\newcommand{\errredA}{9}
\newcommand{\errredL}{12}
\newcommand{\errredLnoetf}{14}
\newcommand{\errredM}{9}
\newcommand{\errredN}{6}
\newcommand{\errredP}{5}
\newcommand{\exitaucmax}{0.82}
\newcommand{\exitaucmin}{0.76}
\newcommand{\exitbasemax}{0.75}
\newcommand{\exitbasemin}{0.65}
\newcommand{\exmqgain}{+3.4}
\newcommand{\exmqhi}{6.7}
\newcommand{\exmqlo}{0.6}
\newcommand{\exmqn}{37}
\newcommand{\ffbuys}{+0.07}
\newcommand{\ffbuyst}{0.87}
\newcommand{\ffsells}{+0.15}
\newcommand{\ffsellst}{2.02}
\newcommand{\ffstocks}{\ensuremath{-}0.04}
\newcommand{\ffstockst}{\ensuremath{-}0.19}
\newcommand{\filedfortyfive}{85}
\newcommand{\filedthirtyone}{0.2}
\newcommand{\filinglagmedian}{45}
\newcommand{\fintextM}{56.7}
\newcommand{\firmcovM}{94}
\newcommand{\firmdmpM}{.03}
\newcommand{\firmdmtM}{2.3}
\newcommand{\firmerrmax}{12.6}
\newcommand{\firmextraM}{+0.47}
\newcommand{\firmextrahiM}{1.00}
\newcommand{\firmextraloM}{\ensuremath{-}0.01}
\newcommand{\firmextramax}{+0.47}
\newcommand{\firmholmM}{.16}
\newcommand{\firmn}{15}
\newcommand{\fundsA}{99}
\newcommand{\fundsL}{57}
\newcommand{\fundsM}{49}
\newcommand{\fundsN}{300}
\newcommand{\fundsP}{59}
\newcommand{\fundsallA}{99}
\newcommand{\fundsallL}{57}
\newcommand{\fundsallM}{49}
\newcommand{\fundsallN}{300}
\newcommand{\fundsallP}{59}
\newcommand{\fundwtgainA}{+2.7}
\newcommand{\fundwtgainL}{+6.4}
\newcommand{\fundwtgainM}{+2.6}
\newcommand{\fundwtgainN}{+2.6}
\newcommand{\fundwtgainP}{+2.3}
\newcommand{\fundwtmedA}{0.0}
\newcommand{\fundwtmedL}{2.7}
\newcommand{\fundwtmedM}{0.0}
\newcommand{\fundwtmedN}{0.3}
\newcommand{\fundwtmedP}{0.7}
\newcommand{\gainA}{+3.5}
\newcommand{\gainL}{+6.2}
\newcommand{\gainLnoetf}{+6.8}
\newcommand{\gainM}{+4.1}
\newcommand{\gainN}{+2.9}
\newcommand{\gainP}{+2.2}
\newcommand{\gainhiA}{5.6}
\newcommand{\gainhiL}{9.5}
\newcommand{\gainhiM}{6.7}
\newcommand{\gainhiN}{4.0}
\newcommand{\gainhiP}{4.2}
\newcommand{\gainloA}{1.6}
\newcommand{\gainloL}{3.0}
\newcommand{\gainloM}{1.8}
\newcommand{\gainloN}{1.9}
\newcommand{\gainloP}{0.5}
\newcommand{\gainlowbaseA}{13.3}
\newcommand{\gainlowbaseL}{12.7}
\newcommand{\gainlowbaseM}{13.8}
\newcommand{\gainlowbaseN}{8.1}
\newcommand{\gainlowbaseP}{6.2}
\newcommand{\gainqhiA}{4.4}
\newcommand{\gainqhiL}{7.4}
\newcommand{\gainqhiM}{4.8}
\newcommand{\gainqhiN}{3.3}
\newcommand{\gainqhiP}{3.0}
\newcommand{\gainqloA}{2.7}
\newcommand{\gainqloL}{5.2}
\newcommand{\gainqloM}{3.3}
\newcommand{\gainqloN}{2.5}
\newcommand{\gainqloP}{1.3}
\newcommand{\gainrestA}{1.2}
\newcommand{\gainrestL}{\ensuremath{-}1.4}
\newcommand{\gainrestM}{1.3}
\newcommand{\gainrestN}{\ensuremath{-}0.3}
\newcommand{\gainrestP}{\ensuremath{-}0.4}
\newcommand{\grattff}{22.0}
\newcommand{\grattseventy}{25.7}
\newcommand{\grattsixty}{23.5}
\newcommand{\grgainff}{+21.6}
\newcommand{\grgainseventy}{+25.1}
\newcommand{\grgainsixty}{+23.0}
\newcommand{\grrecff}{98}
\newcommand{\grrecseventy}{98}
\newcommand{\grrecsixty}{98}
\newcommand{\hbgainA}{+3.2}
\newcommand{\hbgainL}{+5.8}
\newcommand{\hbgainM}{+2.3}
\newcommand{\holdA}{3.2}
\newcommand{\holdL}{22.8}
\newcommand{\holdM}{2.2}
\newcommand{\holdN}{19.9}
\newcommand{\holdP}{17.2}
\newcommand{\holdelliott}{67}
\newcommand{\holdicahn}{68}
\newcommand{\holdpershing}{55}
\newcommand{\holdpooledtestM}{2.5}
\newcommand{\injbasenoisy}{59.0}
\newcommand{\injbaseten}{65.8}
\newcommand{\injmodelnoisy}{88.6}
\newcommand{\injmodelten}{100.0}
\newcommand{\injrecovered}{100}
\newcommand{\injrulenoisy}{88.6}
\newcommand{\labelsA}{402,593}
\newcommand{\labelsL}{195,205}
\newcommand{\labelsM}{305,343}
\newcommand{\labelsN}{465,309}
\newcommand{\labelsP}{227,989}
\newcommand{\lcallA}{+3.5}
\newcommand{\lcallL}{+6.2}
\newcommand{\lcallM}{+4.1}
\newcommand{\lcallhiA}{5.5}
\newcommand{\lcallhiL}{9.5}
\newcommand{\lcallhiM}{6.6}
\newcommand{\lcallloA}{1.6}
\newcommand{\lcallloL}{3.2}
\newcommand{\lcallloM}{1.9}
\newcommand{\lcallquarters}{30}
\newcommand{\lcsixA}{+3.4}
\newcommand{\lcsixL}{+4.8}
\newcommand{\lcsixM}{+2.5}
\newcommand{\lcsixhiA}{6.1}
\newcommand{\lcsixhiL}{8.3}
\newcommand{\lcsixhiM}{4.5}
\newcommand{\lcsixloA}{1.3}
\newcommand{\lcsixloL}{1.4}
\newcommand{\lcsixloM}{0.6}
\newcommand{\lctwelveA}{+3.8}
\newcommand{\lctwelveL}{+5.0}
\newcommand{\lctwelveM}{+3.0}
\newcommand{\lctwelvehiA}{6.2}
\newcommand{\lctwelvehiL}{7.9}
\newcommand{\lctwelvehiM}{5.2}
\newcommand{\lctwelveloA}{1.8}
\newcommand{\lctwelveloL}{2.3}
\newcommand{\lctwelveloM}{1.1}
\newcommand{\lctwentyA}{+3.8}
\newcommand{\lctwentyL}{+4.2}
\newcommand{\lctwentyM}{+4.1}
\newcommand{\lctwentyhiA}{6.3}
\newcommand{\lctwentyhiL}{7.2}
\newcommand{\lctwentyhiM}{6.5}
\newcommand{\lctwentyloA}{1.6}
\newcommand{\lctwentyloL}{1.3}
\newcommand{\lctwentyloM}{2.0}
\newcommand{\lgbA}{62.9}
\newcommand{\lgbL}{55.9}
\newcommand{\lgbM}{60.8}
\newcommand{\lgbN}{56.5}
\newcommand{\lgbP}{56.2}
\newcommand{\lgboracleA}{66.6}
\newcommand{\lgboracleL}{60.7}
\newcommand{\lgboracleM}{63.6}
\newcommand{\lgboracleN}{62.3}
\newcommand{\lgboracleP}{58.6}
\newcommand{\lgboraclegainA}{+7.2}
\newcommand{\lgboraclegainL}{+10.9}
\newcommand{\lgboraclegainM}{+7.0}
\newcommand{\lgboraclegainN}{+8.6}
\newcommand{\lgboraclegainP}{+4.6}
\newcommand{\lgbpooledA}{62.8}
\newcommand{\lgbpooledL}{52.1}
\newcommand{\lgbpooledM}{62.5}
\newcommand{\lgbpooledN}{55.6}
\newcommand{\lgbpooledP}{54.0}
\newcommand{\llmjobs}{30,604}
\newcommand{\llmmemanon}{\ensuremath{-}0.02}
\newcommand{\llmmemgap}{0.023}
\newcommand{\llmmemn}{5,812}
\newcommand{\llmmemnamed}{0.00}
\newcommand{\llmmemp}{.031}
\newcommand{\llmmempos}{9}
\newcommand{\llmmemq}{12}
\newcommand{\llmmemsame}{71}
\newcommand{\llmmemt}{2.5}
\newcommand{\llmunread}{1.1}
\newcommand{\llskillA}{4.4}
\newcommand{\llskillL}{5.9}
\newcommand{\llskillM}{4.1}
\newcommand{\llskillN}{5.1}
\newcommand{\llskillP}{3.6}
\newcommand{\logitM}{58.8}
\newcommand{\lstmM}{61.3}
\newcommand{\majsurprisemaxabst}{1.6}
\newcommand{\medgainA}{0.0}
\newcommand{\medgainL}{2.0}
\newcommand{\medgainM}{0.0}
\newcommand{\medgainN}{0.4}
\newcommand{\medgainP}{0.7}
\newcommand{\medhistM}{7}
\newcommand{\medtrainM}{387}
\newcommand{\mgrhiA}{0.43}
\newcommand{\mgrhiL}{0.18}
\newcommand{\mgrhiM}{0.21}
\newcommand{\mgrloA}{\ensuremath{-}0.12}
\newcommand{\mgrloL}{\ensuremath{-}0.43}
\newcommand{\mgrloM}{\ensuremath{-}0.41}
\newcommand{\mgrnA}{51}
\newcommand{\mgrnL}{34}
\newcommand{\mgrnM}{41}
\newcommand{\mgrrhoA}{+0.16}
\newcommand{\mgrrhoL}{\ensuremath{-}0.14}
\newcommand{\mgrrhoM}{\ensuremath{-}0.12}
\newcommand{\multiquantshareM}{80}
\newcommand{\multisellsalpha}{+0.26}
\newcommand{\multisellst}{2.8}
\newcommand{\naboveA}{22}
\newcommand{\naboveL}{21}
\newcommand{\naboveM}{17}
\newcommand{\naboveN}{114}
\newcommand{\naboveP}{23}
\newcommand{\newA}{18}
\newcommand{\newL}{8}
\newcommand{\newM}{14}
\newcommand{\newN}{7}
\newcommand{\newP}{6}
\newcommand{\newswindowgain}{+2.6}
\newcommand{\newswindowhi}{4.9}
\newcommand{\newswindowlo}{0.8}
\newcommand{\nlowbaseA}{13}
\newcommand{\nlowbaseL}{21}
\newcommand{\nlowbaseM}{11}
\newcommand{\nlowbaseN}{103}
\newcommand{\nlowbaseP}{21}
\newcommand{\nonoverlapbuys}{+0.03}
\newcommand{\nonoverlapsells}{+0.12}
\newcommand{\nopeergainA}{+3.6}
\newcommand{\nopeergainL}{+6.5}
\newcommand{\nopeergainM}{+4.6}
\newcommand{\nopeergainN}{+2.5}
\newcommand{\nopeergainP}{+2.2}
\newcommand{\noverlap}{25}
\newcommand{\nstyletests}{6}
\newcommand{\oraclegainA}{+4.5}
\newcommand{\oraclegainL}{+2.9}
\newcommand{\oraclegainM}{+3.0}
\newcommand{\oraclegainN}{+5.2}
\newcommand{\oraclegainP}{+0.7}
\newcommand{\overlapbuys}{+0.09}
\newcommand{\overlapsells}{+0.16}
\newcommand{\pdbuysA}{\ensuremath{-}0.02}
\newcommand{\pdbuysL}{+0.24}
\newcommand{\pdbuysM}{+0.17}
\newcommand{\pdbuystA}{\ensuremath{-}0.1}
\newcommand{\pdbuystL}{0.9}
\newcommand{\pdbuystM}{0.7}
\newcommand{\pdsellsA}{+0.33}
\newcommand{\pdsellsL}{+0.33}
\newcommand{\pdsellsM}{+0.43}
\newcommand{\pdsellschA}{+0.25}
\newcommand{\pdsellschL}{+0.12}
\newcommand{\pdsellschM}{+0.31}
\newcommand{\pdsellschtA}{1.3}
\newcommand{\pdsellschtL}{0.8}
\newcommand{\pdsellschtM}{1.9}
\newcommand{\pdsellsdlA}{+0.33}
\newcommand{\pdsellsdlL}{+0.29}
\newcommand{\pdsellsdlM}{+0.42}
\newcommand{\pdsellsdltA}{1.4}
\newcommand{\pdsellsdltL}{1.1}
\newcommand{\pdsellsdltM}{3.2}
\newcommand{\pdsellsholmA}{0.532}
\newcommand{\pdsellsholmL}{0.582}
\newcommand{\pdsellsholmM}{0.013}
\newcommand{\pdsellsholmallA}{1.000}
\newcommand{\pdsellsholmallL}{1.000}
\newcommand{\pdsellsholmallM}{0.038}
\newcommand{\pdsellstA}{1.4}
\newcommand{\pdsellstL}{1.4}
\newcommand{\pdsellstM}{3.3}
\newcommand{\pdstocksA}{\ensuremath{-}0.19}
\newcommand{\pdstocksL}{\ensuremath{-}0.56}
\newcommand{\pdstocksM}{\ensuremath{-}0.21}
\newcommand{\pdstockstA}{\ensuremath{-}0.5}
\newcommand{\pdstockstL}{\ensuremath{-}1.2}
\newcommand{\pdstockstM}{\ensuremath{-}0.6}
\newcommand{\pduncondA}{+0.16}
\newcommand{\pduncondL}{\ensuremath{-}0.20}
\newcommand{\pduncondM}{+0.16}
\newcommand{\pduncondtA}{1.5}
\newcommand{\pduncondtL}{\ensuremath{-}1.4}
\newcommand{\pduncondtM}{1.3}
\newcommand{\postbuys}{+0.11}
\newcommand{\postbuyst}{1.43}
\newcommand{\postsells}{+0.04}
\newcommand{\postsellst}{0.61}
\newcommand{\poststocks}{\ensuremath{-}0.05}
\newcommand{\poststockst}{\ensuremath{-}0.24}
\newcommand{\protbaseA}{60.8}
\newcommand{\protbaseL}{50.7}
\newcommand{\protbaseM}{58.1}
\newcommand{\protfundA}{59.7}
\newcommand{\protfundL}{51.1}
\newcommand{\protfundM}{57.0}
\newcommand{\protfundgainA}{\ensuremath{-}1.1}
\newcommand{\protfundgainL}{+0.5}
\newcommand{\protfundgainM}{\ensuremath{-}1.1}
\newcommand{\protlogitM}{56.2}
\newcommand{\protlogitgainM}{\ensuremath{-}1.9}
\newcommand{\protlstmM}{58.8}
\newcommand{\protlstmgainM}{+0.7}
\newcommand{\protpoolA}{63.2}
\newcommand{\protpoolL}{55.0}
\newcommand{\protpoolM}{60.6}
\newcommand{\protpoolgainA}{+2.3}
\newcommand{\protpoolgainL}{+4.4}
\newcommand{\protpoolgainM}{+2.5}
\newcommand{\rawbuys}{+0.07}
\newcommand{\rawbuyst}{0.84}
\newcommand{\rawsells}{+0.21}
\newcommand{\rawsellst}{2.69}
\newcommand{\rawstocks}{\ensuremath{-}0.14}
\newcommand{\rawstockst}{\ensuremath{-}0.78}
\newcommand{\recentM}{52.9}
\newcommand{\recentgainA}{\ensuremath{-}8.6}
\newcommand{\recentgainL}{+1.3}
\newcommand{\recentgainM}{\ensuremath{-}3.7}
\newcommand{\recentgainN}{+1.4}
\newcommand{\recentgainP}{+0.2}
\newcommand{\rhoholdingsA}{+0.26}
\newcommand{\rhoholdingsL}{+0.07}
\newcommand{\rhoholdingsM}{+0.15}
\newcommand{\rhoholdingshiA}{0.48}
\newcommand{\rhoholdingshiL}{0.42}
\newcommand{\rhoholdingshiM}{0.48}
\newcommand{\rhoholdingsloA}{0.03}
\newcommand{\rhoholdingsloL}{\ensuremath{-}0.30}
\newcommand{\rhoholdingsloM}{\ensuremath{-}0.17}
\newcommand{\rqthreemax}{+0.46}
\newcommand{\rqthreemin}{\ensuremath{-}0.05}
\newcommand{\rqthreepos}{21}
\newcommand{\rqthreequartersA}{18}
\newcommand{\rqthreequartersL}{18}
\newcommand{\rqthreequartersM}{18}
\newcommand{\rqthreesig}{3}
\newcommand{\rqtwoall}{\ensuremath{-}0.16}
\newcommand{\rqtwoallhi}{0.29}
\newcommand{\rqtwoalllo}{\ensuremath{-}0.72}
\newcommand{\rqtwochronos}{+0.49}
\newcommand{\rqtwochronoshi}{0.99}
\newcommand{\rqtwochronoslo}{0.08}
\newcommand{\rqtwolatenow}{\ensuremath{-}0.12}
\newcommand{\rqtwolatenowhi}{0.29}
\newcommand{\rqtwolatenowlo}{\ensuremath{-}0.50}
\newcommand{\rqtwolateprev}{\ensuremath{-}0.36}
\newcommand{\rqtwolateprevhi}{0.25}
\newcommand{\rqtwolateprevlo}{\ensuremath{-}0.96}
\newcommand{\rqtwollmnow}{+0.11}
\newcommand{\rqtwollmnowhi}{0.93}
\newcommand{\rqtwollmnowlo}{\ensuremath{-}0.61}
\newcommand{\rqtwollmprev}{\ensuremath{-}0.05}
\newcommand{\rqtwollmprevhi}{0.44}
\newcommand{\rqtwollmprevlo}{\ensuremath{-}0.56}
\newcommand{\rqtwomaxhi}{+0.99}
\newcommand{\rqtwomaxhiabs}{1}
\newcommand{\rqtwonewsnow}{\ensuremath{-}0.50}
\newcommand{\rqtwonewsnowhi}{0.16}
\newcommand{\rqtwonewsnowlo}{\ensuremath{-}1.17}
\newcommand{\rqtwonewsprev}{\ensuremath{-}0.23}
\newcommand{\rqtwonewsprevhi}{0.31}
\newcommand{\rqtwonewsprevlo}{\ensuremath{-}0.79}
\newcommand{\rqtwotimesfm}{+0.13}
\newcommand{\rqtwotimesfmhi}{0.43}
\newcommand{\rqtwotimesfmlo}{\ensuremath{-}0.16}
\newcommand{\ruleM}{49.7}
\newcommand{\shorterrA}{8}
\newcommand{\shorterrL}{9}
\newcommand{\shorterrM}{4}
\newcommand{\shorterrN}{6}
\newcommand{\shorterrP}{5}
\newcommand{\shorterrmax}{9}
\newcommand{\shorterrmin}{4}
\newcommand{\shortgainA}{+3.3}
\newcommand{\shortgainL}{+4.7}
\newcommand{\shortgainM}{+1.8}
\newcommand{\shortgainN}{+2.9}
\newcommand{\shortgainP}{+2.2}
\newcommand{\shortgainhiA}{5.3}
\newcommand{\shortgainhiL}{9.0}
\newcommand{\shortgainhiM}{4.4}
\newcommand{\shortgainhiN}{4.0}
\newcommand{\shortgainhiP}{4.1}
\newcommand{\shortgainloA}{1.5}
\newcommand{\shortgainloL}{0.1}
\newcommand{\shortgainloM}{\ensuremath{-}0.7}
\newcommand{\shortgainloN}{1.9}
\newcommand{\shortgainloP}{0.5}
\newcommand{\smlA}{\ensuremath{-}0.5}
\newcommand{\smlL}{\ensuremath{-}1.9}
\newcommand{\smlM}{\ensuremath{-}2.4}
\newcommand{\smldiffA}{+1.4}
\newcommand{\smldiffM}{\ensuremath{-}0.4}
\newcommand{\smldiffhiA}{5.3}
\newcommand{\smldiffhiM}{3.5}
\newcommand{\smldiffloA}{\ensuremath{-}2.2}
\newcommand{\smldiffloM}{\ensuremath{-}3.8}
\newcommand{\smlhiA}{1.3}
\newcommand{\smlhiL}{1.2}
\newcommand{\smlhiM}{\ensuremath{-}0.9}
\newcommand{\smlloA}{\ensuremath{-}2.0}
\newcommand{\smlloL}{\ensuremath{-}5.4}
\newcommand{\smlloM}{\ensuremath{-}3.9}
\newcommand{\splitlikepct}{0.7}
\newcommand{\stockcihiA}{0.35}
\newcommand{\stockcihiL}{0.20}
\newcommand{\stockcihiM}{0.36}
\newcommand{\stockciloA}{\ensuremath{-}0.35}
\newcommand{\stockciloL}{\ensuremath{-}0.67}
\newcommand{\stockciloM}{\ensuremath{-}0.44}
\newcommand{\stockmdeA}{0.50}
\newcommand{\stockmdeL}{0.62}
\newcommand{\stockmdeM}{0.57}
\newcommand{\styleactgain}{+6.6}
\newcommand{\styleactn}{13}
\newcommand{\stylelsgain}{+2.0}
\newcommand{\stylelsn}{21}
\newcommand{\stylemacrogain}{\ensuremath{-}0.3}
\newcommand{\stylemacron}{3}
\newcommand{\stylemultigain}{+7.9}
\newcommand{\stylemultin}{6}
\newcommand{\stylequantgain}{+4.4}
\newcommand{\stylequantn}{6}
\newcommand{\testrowsA}{99,002}
\newcommand{\testrowsL}{57,486}
\newcommand{\testrowsM}{100,935}
\newcommand{\testrowsN}{286,154}
\newcommand{\testrowsP}{139,892}
\newcommand{\timesfmM}{54.7}
\newcommand{\tqbuysA}{\ensuremath{-}0.81}
\newcommand{\tqbuysL}{\ensuremath{-}0.51}
\newcommand{\tqbuysM}{+0.07}
\newcommand{\tqbuystA}{\ensuremath{-}4.7}
\newcommand{\tqbuystL}{\ensuremath{-}1.3}
\newcommand{\tqbuystM}{0.3}
\newcommand{\tqsellsA}{+0.45}
\newcommand{\tqsellsL}{+1.07}
\newcommand{\tqsellsM}{+0.71}
\newcommand{\tqsellstA}{1.8}
\newcommand{\tqsellstL}{2.9}
\newcommand{\tqsellstM}{3.3}
\newcommand{\trademdehi}{0.26}
\newcommand{\trademdelo}{0.19}
\newcommand{\tunedA}{+3.6}
\newcommand{\tunedL}{+6.2}
\newcommand{\tunedM}{+4.1}
\newcommand{\tunedhiA}{5.6}
\newcommand{\tunedhiL}{9.6}
\newcommand{\tunedhiM}{6.9}
\newcommand{\tunedloA}{1.8}
\newcommand{\tunedloL}{3.3}
\newcommand{\tunedloM}{2.0}
\newcommand{\twowayhiA}{6.3}
\newcommand{\twowayhiL}{9.9}
\newcommand{\twowayhiM}{6.9}
\newcommand{\twowayhiN}{4.1}
\newcommand{\twowayhiP}{4.4}
\newcommand{\twowayloA}{1.4}
\newcommand{\twowayloL}{2.8}
\newcommand{\twowayloM}{1.6}
\newcommand{\twowayloN}{1.7}
\newcommand{\twowayloP}{0.2}
\newcommand{\univholders}{19}
\newcommand{\univnonequity}{2}
\newcommand{\univtype}{63}
\newcommand{\weakabove}{12}
\newcommand{\weakatthi}{8.5}
\newcommand{\weakattlo}{1.2}
\newcommand{\weakdraws}{12}
\newcommand{\weakgainhi}{7.4}
\newcommand{\weakgainlo}{5.1}
\newcommand{\weakotherhi}{3.8}
\newcommand{\weakotherlo}{3.6}
\newcommand{\wfsellsA}{\ensuremath{-}0.35}
\newcommand{\wfsellsL}{+0.43}
\newcommand{\wfsellsM}{+0.50}
\newcommand{\wfsellstA}{\ensuremath{-}0.7}
\newcommand{\wfsellstL}{0.9}
\newcommand{\wfsellstM}{1.0}
\newcommand{\yrdiffM}{+0.7}
\newcommand{\yrerrM}{8}
\newcommand{\yrerrmax}{13}
\newcommand{\yrerrmin}{4}
\newcommand{\yrgainM}{+3.5}
\newcommand{\yrgainhiM}{5.7}
\newcommand{\yrgainloM}{1.6}
\newcommand{\yrpM}{.007}

\newcommand{\how}[1]{\par\smallskip\noindent\textcolor{gray}{\small\emph{How to write it:} #1}\par}

\title{Fact sheets for rewriting the paper\\\large Facts only, no sentences to copy}
\author{}
\date{}

\begin{document}
\maketitle

\noindent\textbf{How to use these.} For each section: (1) read the facts; (2) close
this file and explain the section out loud, as if to a friend; (3) write what you
said, one fact per sentence, using ``we''; (4) check every fact below made it in
and none changed meaning. In the paper, type the number \emph{commands} (shown in
grey boxes where useful), never the numbers themselves.

% ---------------------------------------------------------------
\section*{Abstract (write this last)}
\begin{itemize}
  \item Hook: machine learning predicts 71\% of mutual fund managers' trade
  directions; the trades it cannot predict earn more (Cohen, Lu \& Nguyen).
  \item Question: are hedge funds as predictable?
  \item Data: free SEC filings, 2013--2026; one identical pipeline.
  \item Task: predict buy / hold / sell for each \emph{existing} position among the
  500 most widely held U.S.\ stocks, before the fund's filing is public.
  \item Samples: two hedge fund samples (hand-picked; selected by rule from
  2012--13 Form ADV, \emph{before} the sample starts). Comparisons: active
  long-only managers; N-PORT mutual funds, individually and by family.
  \item Main result: the model fixes only \errredM\% of the errors of each fund's
  majority action (\gainM{} points); the median fund gains nothing.
  \item Checks: the pipeline finds planted signals of moderate strength; more
  training history barely helps.
  \item Comparison: long-only managers and individual mutual funds are not
  significantly more predictable; adding funds up does not significantly change
  this; 2017--2019 is no more predictable than 2022--2026.
  \item Foundation models: three zero-shot models do no better than the baseline
  in per-fund accuracy.
  \item Returns: unpredicted sells by the hand-picked funds fall more before
  disclosure, but this weakens when stocks are matched on size and momentum and
  does not replicate; otherwise no extra return; 18 quarters cannot rule out an
  effect as large as the mutual-fund one.
\end{itemize}

% ---------------------------------------------------------------
\section*{1. Introduction (write second to last)}
\textbf{Paragraph 1: the motivation}
\begin{itemize}
  \item Cohen, Lu \& Nguyen: LSTMs on mutual fund holdings predict 71\% of
  next-quarter trade directions vs 52\% for each fund's most common action;
  the managers and positions the model cannot predict outperform.
  \item Hedge funds are the natural next test: \$5.4 trillion net assets
  (\$12.6 trillion gross) reported to the SEC in early 2025; fewer trading
  constraints; thought to trade on private information.
  \item Their only public record of stock positions is the quarterly 13F, which
  many investors read to copy them.
  \item Time-series foundation models promise forecasts without task-specific
  training, but their value in finance is contested.
\end{itemize}
\textbf{The three questions}
\begin{itemize}
  \item RQ1: how predictable are hedge funds' trades from public information,
  compared with other managers and individual funds measured the same way?
  \item RQ2: do foundation models or news improve the prediction?
  \item RQ3: do the trades that cannot be predicted earn more?
\end{itemize}
\textbf{Framing}: answers are mostly negative; we designed the study so each
negative answer says what size of effect it rules out, and closed the obvious
loopholes (leakage, survivorship, weak models, class imbalance, comparing
across studies).

\noindent\textbf{Contributions (4)}
\begin{enumerate}
  \item First application of the trade-direction design to hedge funds (earlier
  13F work predicts portfolio weights); free public data; five samples:
  \fundsallM{} hand-picked hedge funds, \fundsallA{} Form ADV hedge funds,
  \fundsallL{} long-only managers, \fundsallN{} mutual funds, individually and
  by family.
  \item Hedge funds' trades are hard to predict but not unusual; aggregation
  doesn't change it; the oracle shows a fund's stock choices stay unpredictable
  even when its overall direction is known.
  \item Three zero-shot foundation models (Chronos-2, TimesFM, a Chronos trained
  only on financial data before our test period) and news over several windows.
  \item Positive controls; a multiple-testing-corrected return test with power
  stated; an input-scrambling leakage audit; a classic look-ahead bias that
  inflated our own early results (conditioning on future index membership).
\end{enumerate}

% ---------------------------------------------------------------
\section*{2. Related work}
\textbf{Holdings, herding, skill}
\begin{itemize}
  \item Quarterly holdings are the standard window into what institutions trade.
  \item Herding: Lakonishok--Shleifer--Vishny (1992), Wermers (1999), Sias (2004).
  \item Skill revealed by trades: Chen--Jegadeesh--Wermers (2000); performance
  measured against characteristic-matched benchmarks: Daniel et al.\ (1997).
  \item Hedge funds in 13F: rode the tech bubble and exited before the crash
  (Brunnermeier \& Nagel); only weak evidence they pick better stocks than mutual
  funds (Griffin \& Xu); positions they delay disclosing earn abnormal returns
  (Agarwal et al.; Aragon et al.).
  \item Demand-system models of holdings: Koijen \& Yogo.
  \item Copycats: funds copying disclosed portfolios can match or beat the
  originals (Frank et al.; Verbeek \& Wang).
\end{itemize}
\textbf{Predicting managers' trades}
\begin{itemize}
  \item Cohen et al.: 71\% vs 52\%, Morningstar 1990--2023, per-fund LSTMs;
  unpredictable managers/positions outperform.
  \item Izadifar \& Rahmati: forecast 13F portfolio \emph{weights} for 99 managers
  with temporal graph models; a simple persistence forecast beats all but one.
\end{itemize}
\textbf{Foundation models in finance}
\begin{itemize}
  \item Chronos, Chronos-2, TimesFM: zero-shot forecasting.
  \item Rahimikia et al.: off-the-shelf models forecast daily excess returns
  poorly; training only on financial data works better; per-year models released.
  \item Noguer i Alonso \& Pereira Franklin: on five large U.S.\ stocks, gains over
  a random walk are small and sparse.
  \item Zhang \& Zhang: review LLM stock forecasting; stress pitfalls including data
  leakage, dataset and horizon design, evaluation metrics.
  \item We test foundation models on position sizes, not prices.
\end{itemize}
\textbf{News}: FNSPID headlines scored with FinBERT; policy uncertainty index
(Baker, Bloom \& Davis); Gao et al.\ propose a test for LLM lookahead bias.

% ---------------------------------------------------------------
\section*{3. Data (write this first)}
\textbf{Holdings}
\begin{itemize}
  \item SEC Form 13F data sets: 54 files (quarterly releases to 2023, then
  three-month windows), every report filed July 2013--August 2026, holdings for
  2013Q2--2026Q2.
  \item Who files: managers with at least \$100 million in Section 13(f)
  securities; long positions; within 45 days of quarter-end.
  \item A 13F covers the whole firm, not one fund; we still call each filer a
  ``fund''.
  \item Cleaning: common shares only; latest restatement plus amendments that add
  positions; values converted to dollars (SEC switched from thousands in Jan 2023).
\end{itemize}
\textbf{Samples}
\begin{itemize}
  \item Hand-picked: 49 well-known hedge funds across styles; written rules
  (use the entity that reports holdings, join entity changes, exclude market
  makers / family offices / long-only).
  \item Form ADV: a 2026 list risks favoring survivors, so we pick by rule from
  ADV filings of July 2012--June 2013. Rules: hedge funds $\ge$ \$1B, $\ge$ half
  of private-fund assets, $\ge$ a quarter of regulatory assets; link to 13F by CIK
  or exact name; $\ge$ 10 holdings in June 2013; 100 largest (99 hold universe
  stocks). Follows each adviser's 2013 filer (doesn't follow e.g.\ SAC into Point72).
  \item Long-only: $>75\%$ of assets from mutual funds; essentially no hedge funds;
  $\ge$ \$1B; active share $\ge 0.6$ (excludes Vanguard, Fidelity's firm-wide
  filing). Three ETF sponsors pass (no ETF flag before 2018); results without them
  reported. Same aggregation level as hedge funds.
  \item N-PORT mutual funds: individual funds (the unit of Cohen et al.), from late
  2019; as of Dec 2019, active U.S.\ equity funds ($\ge$ 20 universe stocks, $\ge$
  60\% of net assets, active share $\ge 0.6$, not ETFs), 300 largest; also added up
  to their families (registrants with $\ge$ 2 funds; only complete quarters).
\end{itemize}
\textbf{Stock universe}
\begin{itemize}
  \item 500 largest stocks by total institutional holdings across all 13F filers,
  each quarter; robust to unit errors (drop filers whose price is $>5\times$ off).
  \item Excludes debt, ETFs, funds, ADRs (about \univtype{} of the top 500 per
  quarter) and stocks with $<100$ holders (another \univholders).
  \item Built from SEC data so it keeps later-delisted companies; the S\&P 500
  would lose 10--28\% of early members. Membership fixed at $t-1$.
\end{itemize}
\textbf{Market data and news}: OpenFIGI CUSIP$\to$ticker; Yahoo prices for
3,104 tickers (through June 2026; October 2026 for returns), archived; FNSPID
1.4M headlines for 676 stocks (2013--2023), FinBERT; policy uncertainty index;
Fama--French factors.

% ---------------------------------------------------------------
\section*{4. Method (write this second)}
\begin{itemize}
  \item Labels: compare split-adjusted shares at $t$ vs $t-1$: BUY if up $\ge1\%$,
  SELL if down $\ge1\%$, else HOLD; a full exit is SELL. Same rule as Cohen et al.
  \item Split checks: 9 hand-checked cases pass; at most \splitlikepct\% of cases
  look like a missed split (upper bound); \cusipchangen{} CUSIP-change artefacts.
  \item Task: positions held at $t-1$; features use only filings public by 45 days
  into quarter $t$. New positions excluded (\newM\% of labels), as in Cohen et al.
  \item Split by time: train 2013Q3--2020Q4, validate 2021, test 2022Q1--2026Q2.
  \item Features: last 3 actions, holding age, weight, rank, fund size/turnover/
  buy share/long-run average, stock return, momentum, size, peers' trades (same
  sample), policy uncertainty.
  \item Models: majority, recent majority, repeat, rule; logistic regression;
  LightGBM (settings fixed in advance); LSTM over 8 quarters; Chronos-2, TimesFM,
  finance-trained Chronos (zero-shot, decile forecasts $\to$ score $\to$ label via
  cut-offs tuned on 2021); hybrids with forecasts or news; diagnostics (fund
  weights, no peers, oracle).
  \item Metrics: mean per-fund accuracy (main), pooled accuracy, macro-F1, errors
  fixed, log-loss reduction, BUY-vs-SELL AUC; intervals resample funds.
  \item Leakage audit: scramble all inputs dated $\ge T$, features for $T$ unchanged
  (3 quarters); only 0.03\% of inputs filed too late.
  \item Returns: window from quarter-end to each fund's filing date (max 45 days);
  3 spreads (sells, buys, stocks); 18 quarters; Newey--West, $t$-distribution,
  Holm within and across samples; secondary versions listed in the paper.
\end{itemize}

% ---------------------------------------------------------------
\section*{5. Results}
\subsection*{RQ1}
\begin{itemize}
  \item LightGBM: \baseM\% $\to$ \lgbM\%, gain \gainM{} [\gainloM, \gainhiM].
  LSTM \lstmM\%, logistic \logitM\%. Simple rules worse than baseline.
  \item Median fund gains \medgainM; gain comes from \nlowbaseM{} funds without a
  dominant action (\gainlowbaseM{} vs \gainrestM{} for the rest).
  \item Errors fixed \errredM\%; AUC \aucbsM. Habit, not herding.
  \item Robust: per-fund protocol \protfundgainM; BUY/SELL only \bsgainM; HOLD
  band \hbgainM; without multi-strategy/quant \exmqgain; dollar-weighted \dwgain.
  \item Positive controls: planted rules recovered; weaker planted signals
  \grrecsixty\% recovered; 6 vs all training quarters \lcsixM{} vs \lcallM;
  2017--19 \earlyM; tuning \tunedM.
  \item Comparison: long-only \gainL{} (\errredL\%), N-PORT funds \gainN{}
  (\errredN\%), families \gainP{} (\errredP\%); none significantly different;
  far below 40\%.
  \item Mechanisms: oracle \oraclegainM; with oracle feature \lgboraclegainM;
  opaque-stocks test and number-of-holdings test don't single out hedge funds.
\end{itemize}
\subsection*{RQ2}
\begin{itemize}
  \item Zero-shot per-fund: Chronos-2 \chronosM\%, TimesFM \timesfmM\%,
  finance Chronos \fintextM\% vs baseline \baseM\%. Pooled: Chronos-2
  \chronospooledM\% vs \basepooledM\%. AUC \aucchronosM, \auctimesfmM,
  \aucfintextM{} vs LightGBM \auclgbM.
  \item As features: Chronos-2 \rqtwochronos{} [\rqtwochronoslo, \rqtwochronoshi],
  gone in 2022--23; TimesFM and news (any window) nothing; max upper bound
  \rqtwomaxhi.
\end{itemize}
\subsection*{RQ3}
\begin{itemize}
  \item Hand-picked surprise sells: \pdsellsM\% ($t=\pdsellstM$, Holm
  \pdsellsholmM{} / \pdsellsholmallM); size--momentum matched \pdsellschM;
  within fund \wfsellsM{} ($t=\wfsellstM$); not in Form ADV or long-only.
  \item Link: sells against a fund's direction resemble discretionary trades
  (Alexander, Cici \& Gibson).
  \item Buys, stocks, intensity, unconditional: nothing. Copier window: nothing.
  \item Power: can't rule out the mutual-fund-sized effect.
\end{itemize}

% ---------------------------------------------------------------
\section*{6. Discussion and 7. Conclusion}
\begin{itemize}
  \item Hard to predict beyond habit, but not unusual; not period, not history
  length, not aggregation; remaining difference with Cohen et al.\ is the data.
  \item Practical: holdings disclosures give little predictability beyond habit
  for any manager; unpredicted trades aren't a reliable signal.
  \item Foundation models: short, decision-driven series; next step is
  training on holdings.
  \item Pitfalls: future-index-membership look-ahead; entity changes (Winton).
  \item Threats: what 13F shows; sample selection; new positions excluded;
  data and power.
  \item Future work: LLM reading news (with lookahead check), foundation models
  trained on holdings, predicting new positions.
\end{itemize}

\end{document}
